\chapter{Arquitetura e Design}
\label{chap:arquitetura}

\section{Introdução}
\label{sec:arq-intro}

Este capítulo descreve as decisões arquiteturais e de design tomadas para o desenvolvimento da plataforma FinTwin. São apresentados a arquitetura geral do sistema, o diagrama de componentes, o modelo de dados, o diagrama de classes e as decisões relativas ao design da interface.


\section{Visão Geral da Arquitetura}
\label{sec:arq-visao}

A plataforma segue uma arquitetura de microserviços, dividida em três camadas principais: camada de apresentação (frontend), camada de lógica de negócio (backend/\ac{API}) e camada de dados. A Figura~\ref{fig:arquitetura-componentes} apresenta o diagrama de componentes do sistema.

\begin{figure}[H]
\centering
% Diagrama de Componentes - Arquitetura Gold Lock (corrigido)
\resizebox{\textwidth}{!}{%
\begin{tikzpicture}[
    component/.style={draw, rectangle, rounded corners=3pt, minimum width=3.5cm, minimum height=1.1cm, font=\small, fill=#1!10, draw=#1!60},
    layer/.style={draw, dashed, rounded corners=8pt, inner sep=15pt, fill=#1!5},
    arrow/.style={-{Stealth[length=3mm]}, thick},
    >=Stealth
]

% Camada Frontend
\node[layer=blue, minimum width=12cm, minimum height=5cm] (frontend-layer) at (3, 7) {};
\node[font=\bfseries, anchor=south] at (3, 9.3) {Frontend (React + TypeScript)};

\node[component=blue] (react) at (0.5, 7.8) {React 18 + TS};
\node[component=blue] (tailwind) at (5.5, 7.8) {TailwindCSS};
\node[component=blue] (framer) at (0.5, 6.2) {Framer Motion};
\node[component=blue] (zustand) at (5.5, 6.2) {Zustand (State)};

% Camada Backend
\node[layer=green, minimum width=12cm, minimum height=5cm] (backend-layer) at (3, 1) {};
\node[font=\bfseries, anchor=south] at (3, 3.3) {Backend (Node.js + Express)};

\node[component=green] (api) at (0.5, 1.8) {API REST};
\node[component=green] (auth) at (5.5, 1.8) {Supabase Auth};
\node[component=green] (openbank) at (0.5, 0.2) {Open Banking};
\node[component=green] (notif) at (5.5, 0.2) {Notificações};

% Camada ML (à direita)
\node[layer=orange, minimum width=5cm, minimum height=5cm] (ml-layer) at (12, 1) {};
\node[font=\bfseries, anchor=south] at (12, 3.3) {Serviço ML (Python)};

\node[component=orange] (categorizer) at (12, 1.8) {Categorizador};
\node[component=orange] (irs) at (12, 0.2) {Simulador IRS};

% Serviços Externos (à direita do frontend)
\node[component=purple] (saltedge) at (12, 7.8) {Salt Edge API};
\node[component=purple] (openai) at (12, 6.2) {OpenAI API};

% Camada Dados (em baixo)
\node[layer=red, minimum width=17cm, minimum height=2.5cm] (data-layer) at (6.5, -3.5) {};
\node[font=\bfseries, anchor=south] at (6.5, -2.5) {Camada de Dados};

\node[component=red] (postgres) at (1, -3.8) {PostgreSQL};
\node[component=red] (redis) at (6.5, -3.8) {Redis (Cache)};
\node[component=red] (storage) at (12, -3.8) {Cloud Storage};

% Setas principais
\draw[arrow] (react.south) -- ++(0,-1) -| (api.north);
\draw[arrow] (api.south) -- ++(0,-0.5) -| (postgres.north);
\draw[arrow] (api.south east) -- ++(0,-0.5) -| (redis.north);
\draw[arrow] (openbank.east) -- ++(2,0) |- (saltedge.west);
\draw[arrow] (api.east) -- ++(1,0) |- (categorizer.west);
\draw[arrow] (categorizer.north) -- ++(0,1) -| (openai.south);
\draw[arrow] (categorizer.south) -- ++(0,-1) -| (postgres.north east);

\end{tikzpicture}%
}

\caption{Diagrama de componentes da arquitetura do sistema.}
\label{fig:arquitetura-componentes}
\end{figure}


\section{Stack Tecnológica}
\label{sec:arq-stack}

A Tabela~\ref{tab:stack} apresenta as tecnologias selecionadas para cada camada do sistema, juntamente com a respetiva justificação.

\begin{table}[H]
\centering
\caption{Stack tecnológica do sistema FinTwin.}
\label{tab:stack}
\small
\begin{tabularx}{\textwidth}{l l X}
\toprule
\textbf{Camada} & \textbf{Tecnologia} & \textbf{Justificação} \\
\midrule
Frontend Web & React 18 + TypeScript + TailwindCSS & Ecossistema maduro, tipagem forte, grande comunidade \\
Animações & Framer Motion & Necessário para efeitos Liquid Glass performantes \\
Backend & Node.js + Express + TypeScript & Stack unificada JS/TS, vasto ecossistema fintech \\
Base de Dados & PostgreSQL & Robusto, relacional, ideal para dados financeiros \\
Cache & Redis & Sessões, rate limiting, cache de dados Open Banking \\
IA/ML & Python (scikit-learn) + OpenAI \ac{API} & Classificação + sugestões conversacionais \\
Open Banking & Salt Edge \ac{API} & Cobertura PT, sandbox gratuita, boa documentação \\
Autenticação & Supabase Auth & \ac{OAuth}2, \ac{2FA}, tier gratuito generoso \\
Containerização & Docker + Docker Compose & Ambiente reproduzível, isolamento de serviços, parity dev/prod \\
Deploy & Railway ou Render & Deploy simples, tier gratuito para \ac{MVP} \\
CI/CD & GitHub Actions & Automatização integrada com repositório \\
\bottomrule
\end{tabularx}
\end{table}


\section{Modelo de Dados}
\label{sec:arq-modelo-dados}

O modelo de dados foi desenhado em PostgreSQL, priorizando a integridade referencial e a normalização dos dados financeiros. A Figura~\ref{fig:modelo-er} apresenta o diagrama Entidade-Relacionamento do sistema.

\begin{figure}[H]
\centering
% Diagrama Entidade-Relacionamento - Gold Lock (v3 - sem sobreposições)
\resizebox{\textwidth}{!}{%
\begin{tikzpicture}[
    entitybox/.style={draw, rectangle, rounded corners=2pt, minimum width=4.2cm, fill=blue!5, draw=blue!40, inner sep=6pt, font=\small},
    >=Stealth
]

% === LINHA 1 (topo) ===
% User (esquerda)
\node[entitybox] (user) at (0, 0) {
    \begin{tabular}{l}
    \textbf{User} \\
    \midrule
    \underline{id}: UUID \\
    email: VARCHAR \\
    name: VARCHAR \\
    avatar\_url: TEXT \\
    currency: CHAR(3) \\
    created\_at: TIMESTAMP
    \end{tabular}
};

% BankAccount (centro)
\node[entitybox] (bank) at (8.5, 0) {
    \begin{tabular}{l}
    \textbf{BankAccount} \\
    \midrule
    \underline{id}: UUID \\
    user\_id: UUID (FK) \\
    bank\_name: VARCHAR \\
    iban: VARCHAR \\
    balance: DECIMAL \\
    saltedge\_id: VARCHAR \\
    last\_sync: TIMESTAMP
    \end{tabular}
};

% Transaction (direita)
\node[entitybox] (trans) at (17, 0) {
    \begin{tabular}{l}
    \textbf{Transaction} \\
    \midrule
    \underline{id}: UUID \\
    account\_id: UUID (FK) \\
    category\_id: UUID (FK) \\
    description: TEXT \\
    amount: DECIMAL \\
    date: DATE \\
    is\_recurring: BOOLEAN \\
    ml\_confidence: FLOAT
    \end{tabular}
};

% === LINHA 2 (meio) - com mais espaço vertical ===
% Budget (esquerda)
\node[entitybox] (budget) at (0, -8.5) {
    \begin{tabular}{l}
    \textbf{Budget} \\
    \midrule
    \underline{id}: UUID \\
    user\_id: UUID (FK) \\
    category\_id: UUID (FK) \\
    amount\_limit: DECIMAL \\
    month: DATE \\
    alert\_threshold: INT
    \end{tabular}
};

% SavingsGoal (centro) - afastado do Budget
\node[entitybox] (goal) at (8.5, -8.5) {
    \begin{tabular}{l}
    \textbf{SavingsGoal} \\
    \midrule
    \underline{id}: UUID \\
    user\_id: UUID (FK) \\
    name: VARCHAR \\
    target\_amount: DECIMAL \\
    current\_amount: DECIMAL \\
    deadline: DATE
    \end{tabular}
};

% Category (direita)
\node[entitybox] (cat) at (17, -8.5) {
    \begin{tabular}{l}
    \textbf{Category} \\
    \midrule
    \underline{id}: UUID \\
    name: VARCHAR \\
    icon: VARCHAR \\
    color: VARCHAR \\
    is\_default: BOOLEAN
    \end{tabular}
};

% === LINHA 3 (baixo) ===
% IRSSimulation (centro)
\node[entitybox] (irs) at (8.5, -16) {
    \begin{tabular}{l}
    \textbf{IRSSimulation} \\
    \midrule
    \underline{id}: UUID \\
    user\_id: UUID (FK) \\
    income: DECIMAL \\
    deductions: JSONB \\
    tax\_result: DECIMAL \\
    simulation\_date: DATE
    \end{tabular}
};

% === RELAÇÕES - setas claras sem sobreposição ===

% User -> BankAccount (horizontal, topo)
\draw[->, thick] (user.east) -- node[above, font=\scriptsize, fill=white, inner sep=2pt] {1 : N} (bank.west);

% BankAccount -> Transaction (horizontal, topo)
\draw[->, thick] (bank.east) -- node[above, font=\scriptsize, fill=white, inner sep=2pt] {1 : N} (trans.west);

% Transaction -> Category (vertical, direita)
\draw[->, thick] (trans.south) -- node[right, font=\scriptsize, fill=white, inner sep=2pt] {N : 1} (cat.north);

% User -> Budget (vertical, esquerda) - seta direta pela esquerda
\draw[->, thick] (user.south) -- node[left, font=\scriptsize, fill=white, inner sep=2pt] {1 : N} (budget.north);

% Category -> Budget (horizontal, baixo)
\draw[->, thick] (cat.west) -- node[above, font=\scriptsize, fill=white, inner sep=2pt] {1 : N} (budget.east);

% User -> SavingsGoal - seta pela direita do User, descendo, sem cruzar Budget
\draw[->, thick] (user.south east) -- ++(1.5, 0) -- ++(0, -4.5) -- node[above, font=\scriptsize, fill=white, inner sep=2pt] {1 : N} (goal.north west);

% User -> IRSSimulation - seta pelo centro, descendo
\draw[->, thick] (goal.south) -- node[right, font=\scriptsize, fill=white, inner sep=2pt] {User 1 : N} (irs.north);

\end{tikzpicture}%
}

\caption{Diagrama Entidade-Relacionamento do sistema FinTwin.}
\label{fig:modelo-er}
\end{figure}

As principais entidades do modelo são:

\begin{itemize}
    \item \textbf{User:} Armazena informação do utilizador, incluindo credenciais e preferências.
    \item \textbf{BankAccount:} Representa uma conta bancária ligada via Open Banking, com referência à conexão Salt Edge.
    \item \textbf{Transaction:} Regista cada transação importada, incluindo descrição, montante, data e categoria.
    \item \textbf{Category:} Define as categorias de gastos, incluindo categorias pré-definidas para o mercado português.
    \item \textbf{Budget:} Armazena os orçamentos mensais definidos pelo utilizador por categoria.
    \item \textbf{SavingsGoal:} Representa as metas de poupança com valor alvo, prazo e progresso.
    \item \textbf{IRSSimulation:} Guarda os parâmetros e resultados das simulações de \ac{IRS}.
\end{itemize}


\section{Diagrama de Classes}
\label{sec:arq-classes}

A Figura~\ref{fig:diagrama-classes} apresenta o diagrama de classes simplificado do sistema, focando nas classes principais do domínio.

\begin{figure}[H]
\centering
% Diagrama de Classes Simplificado - Gold Lock (corrigido)
\resizebox{\textwidth}{!}{%
\begin{tikzpicture}[
    classbox/.style={draw, rectangle, rounded corners=2pt, minimum width=4.5cm, fill=#1!8, draw=#1!50, inner sep=6pt, font=\small},
    >=Stealth
]

% Classe User (topo esquerdo)
\node[classbox=blue] (user) at (0, 0) {
    \begin{tabular}{l}
    \textbf{User} \\
    \midrule
    - id: UUID \\
    - email: string \\
    - name: string \\
    - currency: string \\
    \midrule
    + register() \\
    + login() \\
    + updateProfile()
    \end{tabular}
};

% Classe BankAccountService (topo direito)
\node[classbox=green] (bank) at (9, 0) {
    \begin{tabular}{l}
    \textbf{BankAccountService} \\
    \midrule
    - saltEdgeClient: SaltEdge \\
    - accounts: BankAccount[] \\
    \midrule
    + connectAccount() \\
    + syncTransactions() \\
    + getBalance()
    \end{tabular}
};

% Classe TransactionCategorizer (baixo esquerdo)
\node[classbox=orange] (ml) at (0, -7) {
    \begin{tabular}{l}
    \textbf{TransactionCategorizer} \\
    \midrule
    - model: RandomForest \\
    - vectorizer: TfidfVectorizer \\
    - categories: Category[] \\
    \midrule
    + categorize(transaction) \\
    + retrain(corrections) \\
    + getConfidence()
    \end{tabular}
};

% Classe BudgetManager (baixo direito)
\node[classbox=red] (budget) at (9, -7) {
    \begin{tabular}{l}
    \textbf{BudgetManager} \\
    \midrule
    - budgets: Budget[] \\
    - alertThreshold: number \\
    \midrule
    + createBudget() \\
    + checkAlerts() \\
    + getSpending()
    \end{tabular}
};

% Classe IRSSimulator (centro em baixo)
\node[classbox=purple] (irs) at (4.5, -13) {
    \begin{tabular}{l}
    \textbf{IRSSimulator} \\
    \midrule
    - taxBrackets: TaxBracket[] \\
    - deductions: Deduction[] \\
    \midrule
    + simulate(income, deductions) \\
    + compareJointSeparate() \\
    + suggestDeductions()
    \end{tabular}
};

% Relações com labels claros
\draw[->, thick] (user.east) -- node[above, font=\scriptsize, fill=white, inner sep=2pt] {usa} (bank.west);
\draw[->, thick] (bank.south) -- node[right, font=\scriptsize, fill=white, inner sep=2pt] {alimenta} (budget.north);
\draw[->, thick] (bank.south west) -- node[above, font=\scriptsize, fill=white, inner sep=2pt, sloped] {envia para} (ml.north east);
\draw[->, thick] (user.south) -- node[left, font=\scriptsize, fill=white, inner sep=2pt] {configura} (ml.north);
\draw[->, dashed, thick] (ml.east) -- node[above, font=\scriptsize, fill=white, inner sep=2pt] {categorias} (budget.west);
\draw[->, thick] (user.south east) -- node[right, font=\scriptsize, fill=white, inner sep=2pt, pos=0.4] {simula} (irs.north);

\end{tikzpicture}%
}

\caption{Diagrama de classes do sistema FinTwin.}
\label{fig:diagrama-classes}
\end{figure}


\section{Diagrama de Sequência}
\label{sec:arq-sequencia}

A Figura~\ref{fig:sequencia-openbanking} ilustra o diagrama de sequência para o fluxo de ligação de uma conta bancária via Open Banking, que constitui uma das interações mais críticas do sistema.

\begin{figure}[H]
\centering
% Diagrama de Sequência - Ligação Open Banking (corrigido - sem cortes)
\resizebox{\textwidth}{!}{%
\begin{tikzpicture}[
    >=Stealth,
    participant/.style={draw, rectangle, minimum width=2.2cm, minimum height=0.8cm, font=\small\bfseries, fill=#1!10, draw=#1!50},
    msg/.style={->, thick},
    rmsg/.style={->, thick, dashed},
    note/.style={draw, rectangle, rounded corners=2pt, fill=yellow!20, font=\scriptsize, text width=2.8cm, align=center}
]

% Participantes (mais compactos horizontalmente)
\node[participant=blue] (user) at (0, 0) {Utilizador};
\node[participant=green] (frontend) at (3.5, 0) {Frontend};
\node[participant=green] (backend) at (7, 0) {Backend API};
\node[participant=purple] (saltedge) at (10.5, 0) {Salt Edge};
\node[participant=red] (banco) at (14, 0) {Banco};

% Lifelines
\draw[dashed, gray] (0, -0.4) -- (0, -16);
\draw[dashed, gray] (3.5, -0.4) -- (3.5, -16);
\draw[dashed, gray] (7, -0.4) -- (7, -16);
\draw[dashed, gray] (10.5, -0.4) -- (10.5, -16);
\draw[dashed, gray] (14, -0.4) -- (14, -16);

% Mensagens - espaçadas uniformemente
\draw[msg] (0, -1) -- node[above, font=\scriptsize] {1. Seleciona banco} (3.5, -1);
\draw[msg] (3.5, -2) -- node[above, font=\scriptsize] {2. POST /connect} (7, -2);
\draw[msg] (7, -3) -- node[above, font=\scriptsize] {3. Create Connect Session} (10.5, -3);
\draw[rmsg] (10.5, -4) -- node[above, font=\scriptsize] {4. connect\_url} (7, -4);
\draw[rmsg] (7, -5) -- node[above, font=\scriptsize] {5. Redirect URL} (3.5, -5);
\draw[msg] (3.5, -6) -- node[above, font=\scriptsize] {6. Redirect} (0, -6);
\draw[msg] (0, -7) -- node[above, font=\scriptsize] {7. Autenticação no banco} (14, -7);
\draw[rmsg] (14, -8) -- node[above, font=\scriptsize] {8. Consentimento OK} (10.5, -8);
\draw[msg] (10.5, -9) -- node[above, font=\scriptsize] {9. Webhook: success} (7, -9);

% Nota (posicionada sem sobrepor)
\node[note] at (11.5, -10.2) {Backend processa webhook e inicia importação};

\draw[msg] (7, -11.5) -- node[above, font=\scriptsize] {10. GET /accounts} (10.5, -11.5);
\draw[rmsg] (10.5, -12.3) -- node[above, font=\scriptsize] {11. Lista de contas} (7, -12.3);
\draw[msg] (7, -13.1) -- node[above, font=\scriptsize] {12. GET /transactions} (10.5, -13.1);
\draw[rmsg] (10.5, -13.9) -- node[above, font=\scriptsize] {13. Transações (12 meses)} (7, -13.9);
\draw[rmsg] (7, -14.7) -- node[above, font=\scriptsize] {14. Conta ligada} (3.5, -14.7);
\draw[rmsg] (3.5, -15.5) -- node[above, font=\scriptsize] {15. Dashboard atualizado} (0, -15.5);

\end{tikzpicture}%
}

\caption{Diagrama de sequência: ligação de conta bancária via Open Banking.}
\label{fig:sequencia-openbanking}
\end{figure}


\section{Design da Interface --- Liquid Glass}
\label{sec:arq-design}

O design da interface segue o paradigma Liquid Glass, inspirado nas guidelines do iOS 26, adaptado ao contexto de uma aplicação financeira web. Os princípios fundamentais são:

\begin{enumerate}
    \item \textbf{Transparência e profundidade:} Utilização de efeitos de vidro fosco (\textit{backdrop-filter: blur}) em cards e painéis, criando sensação de profundidade e hierarquia visual.
    \item \textbf{Fluidez:} Animações suaves e responsivas implementadas com Framer Motion, priorizando 60 FPS em todas as interações.
    \item \textbf{Acessibilidade:} Cumprimento do nível AA das \ac{WCAG}, com rácio de contraste mínimo de 4.5:1 sobre efeitos de vidro.
    \item \textbf{Modo claro e escuro:} Suporte nativo a ambos os modos, com transição fluída entre eles.
\end{enumerate}

% TODO: Adicionar screenshots/mockups do Figma quando disponíveis


\section{Conclusões}
\label{sec:arq-conclusoes}

A arquitetura proposta garante a separação clara de responsabilidades, a escalabilidade do sistema e a facilidade de manutenção. A escolha de uma stack unificada em TypeScript (frontend e backend) reduz a complexidade do desenvolvimento, enquanto a utilização de Python para o módulo de \ac{ML} permite tirar partido do ecossistema científico da linguagem. O design Liquid Glass diferencia a plataforma visualmente dos concorrentes, contribuindo para uma experiência de utilizador premium.
